\section{Methods}
\label{sec:methods}

\subsection{Pre-registration discipline}

All primary hyperparameters, normalization rules, effect-size thresholds, and the confirmatory dataset family were locked in a pre-registration document (\texttt{PREREGISTRATION.md}) prior to any confirmatory analysis, timestamped and accompanied by a SHA-256 hash recorded at lock time as an internal provenance marker (the hash was computed over the document's content immediately before the hash line itself was appended, so it is not independently recomputable from the committed file as currently stored; it is retained here as a record of the locking procedure followed, not as a reader-verifiable checksum). The pilot cohort (GSE81089) was used to develop the pipeline and is reported separately from the confirmatory family; only the three subsequently-analyzed cohorts (GSE146889, CPTAC-CCRCC, TCGA-LUAD) contribute to the pre-registered, multiple-testing-corrected confirmatory statistics reported in \Cref{sec:results}. The decision rule fixed in advance was that each dataset's result stands on its own merits --- a dataset that fails a gate (e.g., the intrinsic-dimension check below) is reported as blocked or deviated, not silently dropped from the family.

\subsection{Persistent homology pipeline}

For each cohort and omics layer, samples are represented as points in a feature space (genes, methylation probes, or proteins). Two feature spaces are analyzed per cohort/layer: (i) \emph{raw} space, restricted to the top highly-variable features (HVG count fixed at 2000 in the pre-registered configuration; varied in the ablation sweep, \Cref{sec:methods-ablation}); and (ii) \emph{PCA-reduced} space, obtained by projecting the raw feature matrix onto its top 50 principal components (varied in the ablation sweep). A Vietoris--Rips filtration $\Rips(X,\varepsilon)$ is built on the resulting point cloud $X$ in each space, and persistent $H_1$ features are computed; the primary test statistic throughout is $\maxHone$, the longest-lived $H_1$ bar (death $-$ birth) in each filtration. This choice --- a single scalar summary rather than a full persistence-diagram distance (e.g.\ bottleneck or Wasserstein) --- was fixed at pre-registration to keep the statistical comparison against null distributions simple and the effect-size definition unambiguous.

\subsection{Null models}

Two null models are computed in parallel, matched to the same feature space and sample count as the real data at each step:

\begin{description}
\item[Gaussian null.] 500 draws of i.i.d.\ Gaussian noise, dimension- and sample-count-matched to the real data, independently subjected to the identical HVG-selection (where applicable), PCA, and Rips-filtration pipeline as the real data.
\item[Pipeline-symmetric permutation null.] 2000 permutations in which each gene/feature's values are independently shuffled across samples (destroying all cross-feature correlation structure while preserving each feature's marginal distribution exactly), then re-run through the identical HVG-selection, PCA, and filtration pipeline.
\end{description}

The permutation null is the more conservative and more relevant comparison for the $H_1^{\text{stat}}$ (PCA-attribution) test, since it inherits the real data's marginal feature distributions exactly and differs from the real data only in cross-feature dependence structure --- any $\maxHone$ inflation it shows under PCA is attributable purely to the projection step acting on realistic per-feature marginals, not to a mismatched noise model. Both nulls are retained throughout because they answer different questions: the Gaussian null asks whether the signal exceeds what unstructured continuous noise would produce; the permutation null asks whether it exceeds what the same marginal distributions, decorrelated, would produce.

Permutation-count convergence was checked in the ablation sweep (\Cref{sec:methods-ablation}): $z$-scores versus the pipeline-symmetric null were computed at $n_{\text{perm}} \in \{100, 250, 500, 1000, 2000\}$ and found to have converged (bootstrapped 95\% CI stabilized) well before the pre-registered $n=2000$.

\subsection{Effect size and significance thresholds}

Effect size against each null is reported as a $z$-score: $(\text{observed} - \text{null mean})/\text{null SD}$. The pre-registered significance threshold is $z>3.0$ in every comparison (both null models, both feature spaces), chosen to clear the pilot's weakest single margin with room to spare while remaining well above conventional large-effect benchmarks. For the confirmatory family (GSE146889, CPTAC-CCRCC, TCGA-LUAD), the resulting $p$-values across up to 20 planned tests are corrected for multiple comparisons via Benjamini--Hochberg FDR control at $\alpha=0.05$~\cite{benjaminihochberg1995}; a test is reported as passing $H_0^{\text{stat}}$ only if it clears both the raw $z>3.0$ threshold and the BH-adjusted significance criterion.

\subsection{Intrinsic dimension and Johnson--Lindenstrauss regime gate}

Before PCA is applied, the intrinsic dimensionality of each feature space is estimated using two independent estimators --- the Levina--Bickel maximum-likelihood estimator~\cite{levinabickel2004} and the TwoNN estimator~\cite{facco2017} --- and compared against the ambient (post-HVG) dimension. When the ratio $d_{\text{int}}/d_{\text{amb}} \geq 0.3$, the data is judged to sit outside the regime in which Euclidean Vietoris--Rips distances behave predictably under linear projection (the Johnson--Lindenstrauss heuristic~\cite{johnsonlindenstrauss1984} is not expected to hold), and a spectral-cascade metric --- in the spirit of recent spectral-distance approaches to persistent homology on high-dimensional, low-intrinsic-dimension data~\cite{damrich2024} --- is substituted for the naive Euclidean distance in that space's filtration. This gate was triggered for TCGA-LUAD's methylation PCA space ($d_{\text{int}}/d_{\text{amb}} = 0.320$; \Cref{sec:results}), where it is flagged as a pre-registered deviation rather than an ad hoc fix.

\subsection{Datasets}

Five cohorts, spanning four tissue/tumor types and three omics modalities, were analyzed:

\begin{itemize}
\item \textbf{GSE81089} (pilot, not part of the confirmatory family): NSCLC bulk RNA-seq, $n=218$ (199 tumor / 19 normal).
\item \textbf{GSE146889}: colorectal/endometrial/ovarian MMR-deficiency cohort, bulk RNA-seq, $n=176$ (91 tumor / 85 paired-normal).
\item \textbf{CPTAC-CCRCC} (PDC000127): renal clear-cell carcinoma, mass-spectrometry proteomics, $n=194$ (110 tumor / 84 normal), no substitution required.
\item \textbf{TCGA-LUAD}: lung adenocarcinoma, two omics layers --- RNA-seq ($n=116$) and DNA methylation ($n=36$ paired samples, 58 cases with both layers available).
\end{itemize}

Normalization follows modality-specific pre-registered rules: bulk RNA-seq is $\log_1p$-transformed (FPKM or TPM as available per cohort); DNA methylation is logit-transformed on the $\beta$-value scale (clipped to $[10^{-6}, 1-10^{-6}]$ to avoid boundary singularities); CPTAC proteomics is median-normalized then $\log_2$-transformed. Missing values are zero-filled in the pre-registered primary analysis; three alternative imputation rules (mean, median, drop) are compared in the ablation sweep (\Cref{sec:methods-ablation}).

TCGA-LUAD required four declared deviations from the pre-registered protocol, each pre-specified as a contingency rather than an ad hoc decision after seeing results: (D1) the ``bottom-variance'' negative-control gene set is selected among genes with nonzero variance only (5{,}556 of 60{,}660 genes have exactly zero variance and would trivially fail any test); (D2) methylation $\beta$-values are clipped before the logit transform; (D3) methylation PCA is capped at 35 components rather than the pre-registered 50, since only 35 are available given the sample count, and is treated as a rotation-only isometry check rather than an independent significance test; (D4) the spectral-cascade metric is applied to the methylation PCA space per the intrinsic-dimension gate above.

\subsection{Ablation / hyperparameter-sensitivity sweep}
\label{sec:methods-ablation}

To establish the boundary of validity around the pre-registered operating point (HVG$=2000$, PC$=50$), an 18-configuration sweep was run on the pilot cohort (GSE81089), varying: HVG gene count ($\{500, 1000, 2000, 4000\}$, PC count fixed at 50); PCA component count ($\{10, 25, 50, 100\}$, HVG fixed at 2000); missing-value imputation rule (zero, mean, median, drop); and a $2\times2$ gene-count $\times$ PC-count interaction check at the corners of the grid (HVG$\in\{500,4000\}$ $\times$ PC$\in\{10,100\}$) to test whether the two hyperparameters' effects are additive or interact superadditively. Permutation count was reduced to $n=500$ in the sweep (from the pre-registered $n=2000$) to keep the 18-configuration sweep computationally tractable; convergence at this reduced count was separately verified (above).

\subsection{Confound-attribution audit protocol}
\label{sec:methods-confound}

Where an $H_1^{\text{stat}}$ signal is found to survive both null models, a separate question is whether that signal reflects genuine topological structure in the omics manifold, or is instead a re-description of the tumor/normal class-mean separation --- a shift that a simple linear classifier can already recover, in every cohort tested here, at or near ceiling cross-validated AUC. The audit protocol, applied here to all four cohorts (\Cref{sec:results-confound}, \Cref{sec:results-confound-generalization}), consists of four complementary checks:

\begin{enumerate}
\item \textbf{Within-class decomposition.} The $\maxHone$ statistic is recomputed on the tumor-only subpopulation alone (excluding normal samples entirely), and separately on the normal-only subpopulation. If the mixed-population signal is driven purely by the tumor/normal class-mean gap, it should collapse when either subpopulation is analyzed in isolation.
\item \textbf{Within-stratum control.} The tumor-only subpopulation is further stratified into quartiles along the direction of maximal tumor/normal class-mean separation (the top discriminant axis, e.g.\ PC1 where correlated with class), and the $H_1^{\text{stat}}$ test is repeated within each quartile. A signal that survives uniformly across quartiles cannot be attributed to residual variation along that specific confound direction.
\item \textbf{Residualization control.} The tumor/normal class-mean-shift direction is estimated and projected out of the feature matrix (residualization), and the $H_1^{\text{stat}}$ test is repeated on the residualized data. A 5-fold cross-validated logistic regression classifier trained to recover tumor/normal status from the residualized data is also evaluated; its cross-validated AUC collapsing toward chance is evidence the confound direction has been successfully removed without also removing the topological signal. We note that linear confound-regression procedures of this kind have themselves been shown, in other machine-learning contexts, to introduce leakage or distort downstream nonlinear statistics rather than cleanly removing the confound~\cite{hamdan2023}; we do not use a nonlinear downstream model here (only the linear residual's own topology and a linear classifier check), which limits but does not eliminate this concern (\Cref{sec:discussion}).
\item \textbf{Block-bootstrap confidence interval.} A block-bootstrap procedure (resampling at the sample level, preserving each feature's within-sample structure) is used to construct a 95\% confidence interval on the residual signal-beyond-null quantity after residualization; an interval excluding zero is evidence the residual signal is not a bootstrap artifact of the small effective sample size in some strata.
\end{enumerate}

This protocol was applied identically across all four cohorts (\Cref{sec:results-confound}, \Cref{sec:results-confound-generalization}). GSE146889's own replication report had already surfaced a specific structural reason its result might diverge from the pilot's: its dominant $H_1$ loop is significantly enriched for tumor-status samples (Fisher's exact $p=2.9\times10^{-9}$), unlike the pilot's small, non-enriched loop; the audit confirms this concern is real, though only partially, rather than resolving it in either direction.
